\documentclass[11pt]{article}
\usepackage[margin=1in]{geometry}
\usepackage{setspace}
\usepackage{parskip}
\usepackage{amsmath,amssymb}

\begin{document}

\begin{center}
{\Large Note on Regulatory Changes to Auditor-Change Disclosures (2001--2023)}
\end{center}

The short answer is yes: there was a meaningful regulatory change during the sample period, but it is probably not a clean policy shock for this paper's purposes.

The key institutional change is the SEC's 2004 overhaul of Form 8-K in Release No. 33-8400, effective August 23, 2004. That reform expanded Form 8-K generally and accelerated the filing deadline for most items to four business days. However, for auditor changes specifically, the SEC stated that Item 4.01 was ``substantively the same as former Item 4,'' with only limited textual cleanup and reorganization. Thus, the 2004 reform changed the timing and broader architecture of Form 8-K much more than it changed the substantive content requirements governing auditor-change disclosures.

Accordingly, the answer to the identification question is nuanced.

First, the regulatory environment did change during 2001--2023, so it is reasonable to acknowledge that the disclosure regime was not fully static over the sample period.

Second, the most important 2004 change appears to be a filing-timeliness shock rather than a sharp change in the informational content of Item 4.01 itself. That distinction matters. A pre/post-2004 comparison would not cleanly isolate a change in the substance of auditor-change disclosure requirements; instead, it would bundle together faster filing, the post-SOX regulatory environment, early PCAOB-era developments, and broader shifts in investor processing of 8-K disclosures.

Third, for that reason, I would not recommend presenting the 2004 reform as the paper's central policy shock or as a clean quasi-experiment. A referee is likely to object that the treatment is effectively universal and that the SEC itself described the auditor-change item as substantively unchanged. A design centered on this break would therefore have weak attribution.

That said, the reform is still useful empirically as a secondary institutional check. The paper can exploit the 2004 break in three modest ways.

\begin{enumerate}
    \item Include a \textit{Post-2004} indicator as a control, to absorb any level shift associated with the accelerated filing regime.
    \item Estimate a \textit{Polarization $\times$ Post-2004} interaction as a supplementary test. If the polarization effect strengthens after auditor-change filings become more timely, that is an interesting institutional amplification result.
    \item Split the sample into pre-2004 and post-2004 periods to show that the main result is not driven solely by the later disclosure regime.
\end{enumerate}

The framing, however, should remain disciplined. The 2004 change is best described as a regulatory-timing shift, not as a sharp exogenous shock to the substantive content of auditor-change disclosures.

A clean way to state this in the paper would be:

\begin{quote}
Although the SEC accelerated Form 8-K filing in 2004, Item 4.01 itself remained substantively similar to the prior accountant-change disclosure requirement. Accordingly, our design does not rely on the 2004 reform as a sharp policy shock; instead, we use pre/post-2004 analyses as institutional-stability checks to verify that the main polarization result is not driven solely by the later reporting regime.
\end{quote}

In short, yes, there was a regulatory change during the sample period, and yes, it can be used as a robustness or institutional-regime check. But no, it is probably not strong enough to serve as the paper's main policy-shock identification strategy.

\end{document}
